
\begin{center}
{\color{accent}\Huge\bfseries PHYS4055}\\[3mm]
{\LARGE Atomic Physics and Quantum Optics}\\[2mm]
{\large Open-book examination: four-part course overview}\\[2mm]
{\small Source-reviewed revision edition \quad 13 November 2026}
\end{center}
\vspace{4mm}
\begin{methodbox}
\textbf{Purpose.} A compact conceptual reference: each topic explains
the physical system, its governing equations, the assumptions and how
to use the equations in calculations. Use this with the separate
worked-example and past-exam books. It is not a substitute for the
original slides when a derivation depends on a figure or convention.
\end{methodbox}
\section*{The four parts of the course}
\begin{tabularx}{\linewidth}{@{}p{.20\linewidth}Xr@{}}
\toprule
Course part & Main questions & Pages\\\midrule
\textbf{I. Quantum optics} & What is a quantum field mode? How do Fock/coherent,
thermal and squeezed states differ? How do beam splitters, photodetectors,
quadratures and quasiprobabilities describe optical experiments? & 2--4\\[2mm]
\textbf{II. Applications of optics} & How do photonic qubits and entanglement
enable quantum communication? What do open-system noise, conditional measurements,
atom--cavity coupling and quantum metrology accomplish? & 5--9\\[2mm]
\textbf{III. Atomic physics} & How do we calculate atomic orbitals and energies?
How do angular-momentum addition, exchange symmetry, fine/hyperfine interactions
and external fields determine the spectrum? & 10--12\\[2mm]
\textbf{IV. Applications of atomic physics} & Why do spectral lines have
selection rules and lifetimes? How can lasers cool atoms? How do atoms test
fundamental symmetries and operate as extremely accurate clocks? & 13--15\\\bottomrule
\end{tabularx}
\section*{How to use this booklet during the exam}
\begin{enumerate}
\item \textbf{Choose a physical model.} Number-state optics, Gaussian optics,
a quantum channel, a two-level atom, or a coupled angular-momentum multiplet
have different natural bases.
\item \textbf{Write conventions.} For optics we use
\(x=a+a^\dagger,\ p=i(a^\dagger-a)\), so vacuum variance is 1.
For atomic physics state \(l,s,L,S,J,I,F\) and the basis.
\item \textbf{Start from the boxed equation and the method box.}
A named model alone is not an answer: show enough algebra to connect
the question's given quantities to the requested observable.
\item \textbf{Perform a physical check:} conservation of probability,
state normalisation, commutators, degeneracy sums, operator parity,
triangle rules, limiting cases, dimensions and units.
\end{enumerate}
\section*{Source audit and notation}
The underlying lecture slides were reviewed across
Weeks 1--6 (optics and applications), Weeks 8--9 (atomic foundations),
Week 10 (spectroscopy/cooling), Week 12 (symmetry tests) and Week 13
(clocks), alongside the available lecture transcripts.
Individual pages cite slide packs and page ranges. Slide/text extraction
can misread equations, particularly in annotated material: the
original PDFs remain authoritative. The **source map** next to this
booklet records coverage and known issues.
\vfill
{\small\color{gray}
This overview is educational revision material, not an official
formula sheet or a prediction of exam questions.
Date of preparation: 13 November 2026.}
